problem:  
Let \((X,d,\mu)\) be a compact geodesic metric measure space satisfying the **Bishop–Gromov volume comparison** with parameters \((k,N)\) globally; namely, for every \(x\in X\),  
\[
r \mapsto \frac{\mu(B(x,r))}{v_{k,N}(r)}
\]
is nonincreasing for all \(r>0\).  
Assume that for every \(x\in X\), all tangent cones \(\mathrm{Tan}(X,x)\) (in the pointed measured Gromov–Hausdorff sense) exist and are **metric cones** \(C(Z_x)\) over compact length spaces \(Z_x\). Denote by \(\dim_H(Z_x)\) the Hausdorff dimension of the link \(Z_x\) with respect to its natural induced metric.  

Open problem:  
Does the Bishop–Gromov volume comparison property alone imply that  
\[
\dim_H(Z_x) = \text{constant for }\mu\text{-almost every }x\in X ?
\]  
Equivalently, can one construct a compact metric measure space satisfying exact Bishop–Gromov volume monotonicity but having tangent cones of metric‑cone type whose links \(Z_x\) possess **nonconstant Hausdorff dimensions** on a positive‑measure set of points?  

Why is it a "good" problem:  
This problem clarifies precisely how far the global Bishop–Gromov comparison principle constrains the local infinitesimal geometry of singular spaces.  

- **Profound insight:** It isolates a core issue: whether macroscopic curvature control via volume ratios can enforce *microscopic dimension rigidity* of tangent structures, without invoking stronger analytic curvature–dimension conditions. A positive answer would show that volume comparison alone governs not just growth but also the local dimensional structure of space; a counterexample would delineate the sharp boundary of this principle.  
- **Bridge between fields:** It integrates comparison geometry, geometric measure theory, and Gromov–Hausdorff analysis, forcing interaction between global volume inequalities and local dimension theory. Any progress would push the understanding of how analytic curvature notions emerge from purely geometric comparison data.  
- **Extreme simplicity:** The question can be stated with elementary geometric notions—balls, measures, limits, and Hausdorff dimension—yet its resolution would likely require new structural tools at the interface of metric geometry and analysis on singular spaces.